\scanpage{37}\mspage{36}
as soon as $K$ is a final object, in the category of all commutative squares
\diagram{\begin{tikzcd}[ampersand replacement=\&,column sep=large]
\unsure{{+}J'|J'}\arrow[r,"i'"]\arrow[d,"h_J"] \& I'\arrow[d,"h"]\\
J\arrow[r] \& I
\end{tikzcd}}
\deleted{which in} \inserted{of ordered sets (or preordered)} which satisfy the preceding condition. That should imply $(b,c,d)\Longleftrightarrow a)$.

\uline{18 June}\quad In all these questions, the point of view of ultr. topologies \inserted{(i.e. \uuline{pre}ordered topologies)} is perhaps the most practical, without always taking the trouble to spell out laboriously in terms of ordered index sets. A ``figure'' \inserted{of $L$} is made explicit as a pair $(\deleted{A}S,T)$ of a subset $S$ of $L$, and an ultr.topology $T$ on $S$ (topology where \uline{every} union of closed sets is closed, i.e \uline{every} intersection of open sets is open, i.e. the intersection of the neighborhoods \inserted{[\unsure{open, if one wishes|open if one wishes}]} of every point is a neighborhood, or again that every point has a smallest neighborhood (necessarily open). If $(S,T),(S',T')$ are like that, to say that they are compatible means that $T,T'$ induce the same topology on $S\cap S'$, \inserted{which is closed in $S'$ (for $T$) and in $S'$ (for $T'$)} and then there's a unique
